% Configuration des marges
\usepackage[margin=0.8in]{geometry} % Définit les marges du document

% Personnalisation des titres de section et table des matières
\usepackage{sectsty, tocloft, hyperref} % Importe les packages nécessaires pour personnaliser les titres et la table des matières
\hypersetup{
    pdfauthor={Kevin Delfour}, % Définit l'auteur du PDF
    pdftitle={L'envers du décor des Entreprises libérée}, % Définit le titre du PDF
    pdfkeywords={entreprise, libérées, utopie, envers, décor}, % Définit les mots-clés du PDF
    hidelinks % Cache les liens dans le PDF
}

% Cadres, tailles de police et en-têtes/pieds de page personnalisés
\usepackage{framed, anyfontsize, fancyhdr} % Importe les packages nécessaires pour personnaliser les cadres, les tailles de police et les en-têtes/pieds de page

% Police Garamond et paramètres de langue
\usepackage{ebgaramond, float} % Importe la police Garamond et le package float pour la gestion des flottants
\usepackage[french]{babel} % Définit le français comme langue principale
\usepackage[T1]{fontenc} % Utilise la codification de caractères T1

% Couleurs personnalisées et graphiques TikZ
\usepackage[dvipsnames]{xcolor} % Importe le package xcolor avec les noms de couleurs dvips
\usepackage{tikz, microtype, GS1, qrcode} % Importe les packages nécessaires pour les graphiques TikZ, la microtypographie, les codes-barres GS1 et les QR codes
\usetikzlibrary{calc} % Utilise la bibliothèque calc de TikZ pour les calculs de coordonnées

\usepackage[labelformat=empty]{caption} % Importe le package caption avec un format de label vide
\captionsetup{justification=raggedleft, font={tiny, color=gray}, singlelinecheck=false} % Configure les légendes des figures

% Contenu : texte aléatoire, boucles, images et commandes supplémentaires
\usepackage{lipsum, pgffor, graphicx, etoolbox} % Importe les packages nécessaires pour le contenu : texte aléatoire, boucles, images et commandes supplémentaires

% Configuration de la mise en page
\renewcommand{\familydefault}{\sfdefault} % Utilise la police sans serif par défaut
\setlength{\parskip}{\baselineskip} % Définit l'espacement entre les paragraphes
\setlength{\parindent}{0pt} % Supprime l'indentation des paragraphes

% Configuration de la table des matières
\setcounter{tocdepth}{2} % Définit la profondeur de la table des matières
\renewcommand{\cftpartfont}{\normalfont\bfseries\Large} % Configure la police des parties dans la table des matières
\renewcommand{\cftchapfont}{\normalfont\bfseries} % Configure la police des chapitres dans la table des matières
\renewcommand{\cftsecfont}{\normalfont} % Configure la police des sections dans la table des matières
\renewcommand{\cftsubsecfont}{\normalfont} % Configure la police des sous-sections dans la table des matières
\renewcommand{\cftsubsubsecfont}{\normalfont\small} % Configure la police des sous-sous-sections dans la table des matières

% Configuration des en-têtes et pieds de page personnalisés
\pagestyle{fancy} % Utilise le style de page fancy
\fancyhead{} % Supprime l'en-tête par défaut
\fancyhead[LO, L]{ \color{gray}\scriptsize\leftmark} % Configure l'en-tête gauche
\fancyfoot{} % Supprime le pied de page par défaut
\fancyfoot[R]{\scriptsize\thepage} % Configure le pied de page droit
\renewcommand{\headrulewidth}{0pt} % Supprime la ligne sous l'en-tête

\newcommand{\listfootnotesname}{Liste des sources}% 'List of Footnotes' title
\newlistof[chapter]{footnotes}{fnt}{\listfootnotesname}% New 'List of...' for footnotes

\newcommand{\insertImage}[2]{
    \markboth{}{} % Supprime les marques de tête
    \begin{figure}[H]
        \center % Centre l'image
        \includegraphics[keepaspectratio, width=\textwidth, height=\textheight]{images/#1} % Insère l'image
        \captionsetup{justification=raggedright, singlelinecheck=false} % Aligne la légende sur la gauche
        \caption{#2} % Ajoute une légende à l'image
    \end{figure}
    \newpage % Crée une nouvelle page après l'image
}

\makeatletter
\renewcommand{\thefigure}{\@arabic\c@figure}
\makeatother


\let\oldfootnote\footnote % Save the old \footnote{...} command
\renewcommand\footnote[1]{
    \oldfootnote{#1}% Make a regular footnote with added margin.
    \addcontentsline{fnt}{footnotes}{\protect
        \numberline{\thepart.\thefootnotes}\hspace{1,8em}#1}% Add the 'List of...' entry with a margin between the number and the note, and include the part number in the numbering.
}

\makeatletter
\@addtoreset{chapter}{part} % Réinitialise le compteur de chapitres à chaque nouvelle partie
\makeatother